\documentclass[UTF8]{ctexart}
\usepackage{geometry}
\usepackage{hyperref}
\usepackage{listings}
\usepackage{xcolor}

\geometry{a4paper, left=2.5cm, right=2.5cm, top=2.5cm, bottom=2.5cm}

\title{Agent开发周报 260415}
\author{刘德辰}
\date{2026年4月15日}

\begin{document}

\maketitle

\section{报告开发基本对齐到车辆完成状态：驾驶员已基本完成，单位/线路/事故补齐到同类状态}

\textbf{背景}：上周已把多类主画像报告切到 MCP，并统一数据源与分区策略；本周的重点是把\textbf{驾驶员、单位、线路、事故}四类报告整体推进到与现有\textbf{车辆报告}接近的完成状态，做到可以开始系统测试，使它们在「取数—生成—展示—多步承接」上的实现方式尽量一致，减少后续维护时每类报告各走一套的分裂状态。这样也能避免架构和管线状态不一致，更方便后续分工测试和迭代。

\textbf{方案说明}：本周围绕「对齐车辆报告完成状态」推进了四条线：一是完成驾驶员画像报告生产侧收尾；二是补齐单位管理类报告；三是让线路报告技能、数据源与行为约束向其它报告靠齐；四是在周中继续把事故调查报告的提示词、数据源说明、结构化归一化与运行时接入补齐到与车辆/驾驶员同级别的完成状态，并统一接入结构化输出 formatter，使不同报告类型在同一套框架下演进。

\textbf{完成情况说明}：
\begin{itemize}
\item \textbf{驾驶员报告}：已经基本对齐到车辆报告的完成状态，报告主体开发可以认为已经完成，当前可以开始系统测试。剩余主要是路由召回上的继续回归、真实会话观察与迭代。
\item \textbf{单位报告}：主体管线已经补齐到和车辆/驾驶员同类的状态，可以进入并行测试；但多轮补参和上下文承接还有尾项，需要继续修路由召回细节。
\item \textbf{线路报告}：整体能力已经基本对齐到统一框架，取数、生成和展示方式都在向车辆报告靠齐；当前更像是「开发基本到位、系统验证还需要补齐」，还需要补一轮更系统的召回与多轮场景测试。
\item \textbf{事故报告}：这周也已补齐到与车辆/驾驶员接近的完成状态，报告主体开发基本到位，可以按同类报告的方式开始测试和迭代。剩余更偏路由验证、边缘场景回归和线上观察。
\item \textbf{统一支撑层}：各类结构化报告继续共用统一的输出格式化与展示整形能力，进一步降低了不同报告之间的实现差异，便于后续分工测试和并行迭代。
\end{itemize}

\section{路由与查数容错：结构化 lookup 失败回路由、车辆/驾驶员路由规则与多步续写}

\textbf{背景}：结构化查数或解析失败如果只停留在 worker 内部处理，系统很容易和用户当前真正想做的事脱节，例如用户已经在换对象、补参数、改问法，或者从「生成报告」切到「先咨询一下」。同时，多步开场/退场接入后，报告类对话在上下文承接、意图切换和续写判断上也更容易暴露路由问题。

\textbf{方案说明}：本周继续把这类「本该由 router 重新判断」的情况收回到 router 统一处理，让系统在多轮对话里更像是在跟着用户意图走，而不是卡在上一步的工具或 worker 里；同时继续收紧车辆/驾驶员报告与车辆专家相关的路由边界，减少误分流。

\textbf{功能说明}：
\begin{itemize}
\item \textbf{解决了查数失败后容易卡死的问题}：当结构化 lookup 没命中、解析失败或拿不到有效结果时，不再默认停留在原 worker 里硬撑，而是回到 router 重新判断，让系统有机会根据用户当前上下文改走澄清、咨询或别的报告路径。
\item \textbf{解决了多轮对话里上下文承接不稳的问题}：围绕多步报告 continuation 做了细化，使系统在「上一轮已经在生成/解释某个报告，这一轮继续追问、补参、纠错或续写」时，更容易延续正确上下文，而不是把会话当成一条全新请求。
\item \textbf{解决了意图转换时容易误分流的问题}：当用户在同一会话里从报告切到咨询、从一个对象切到另一个对象，或从泛问切到明确报告需求时，router 的约束比之前更清楚，降低了「还沿着旧意图继续走」的风险。
\item \textbf{继续收紧了车辆/驾驶员报告与专家路径的边界}：这周一起调整了驾驶员报告相关 prompt、车辆报告路由规则、车辆专家判定逻辑，以及「单车风险/画像类查询应优先走专家而不是直接出正式报告」的分流规则，让报告、专家、咨询三类路径的职责更清楚。
\end{itemize}

\section{Runtime 拆分加速：HTTP、Query/OpenAI、Auth、KB、Rule、ToolProvider 与 prompt 模块外提}

\textbf{背景}：本周在 sessions/AB 等已迁出的基础上，把 \textbf{HTTP 路由、查数与 OpenAI 调用、鉴权、KB、规则域、工具 provider} 等大块从 Runtime 中剥离，并拆分 runtime prompt 模块与文档，降低后续改路由与改工具链时的牵连面。

\textbf{方案说明}：按职责拆模块，让运行时的组合逻辑与提示词资产之间的边界更清晰。更方便迭代运行时 prompts 与工具链。

\textbf{功能说明}：
\begin{itemize}
\item 从 runtime 拆出 HTTP router（\texttt{split http router from runtime}）
\item 从 runtime 拆出 query data 与 openai（\texttt{split query data and openai from runtime}）
\item 从 runtime 拆出 auth（\texttt{split auth from runtime}）
\item 从 runtime 拆出 kb（\texttt{split kb from runtime}）
\item 从 runtime 拆出 rule domain（\texttt{split rule domain from runtime}）
\item 从 runtime 拆出 tool provider（\texttt{split tool provider from runtime}）
\item 拆分 runtime prompt 模块与文档（\texttt{split runtime prompt modules and docs}），并增加 prompts 引用说明（\texttt{add prompts ref}）
\end{itemize}

\section{专家能力扩展与元数据：新增驾驶员专家、KB 工具信息进入 metadata}

\textbf{背景}：车辆专家路径已运行多轮，本周将 \textbf{驾驶员专家} 纳入与 router/runtime 的衔接；同时希望在元数据中显式暴露 KB 相关工具信息，便于上层路由与观测。

\textbf{方案说明}：新增 driver expert 能力并与现有专家体系并列；在 metadata 中补充 kb tools 信息。

\textbf{功能说明}：
\begin{itemize}
\item 新增驾驶员专家（\texttt{add driver expert}）
\item 在 metadata 中增加 KB 工具信息（\texttt{add kb tools info to metadata}）
\end{itemize}

\section{报告前端呈现与 MCP 增强：Markdown 检测与样式、指标来源、车辆建议与分区参数}

\textbf{背景}：结构化报告混排 Markdown 时，需要可靠的检测与标题/表格样式；指标「看起来像可点」与实际来源需一致；车辆报告侧希望利用 MCP 建议能力 enrich，同时约束指标来源避免误引用。

\textbf{方案说明}：修复 Markdown 报告渲染识别，补充 heading/table 样式；改进 report indicator 的展示与来源约束；为车辆报告接入 MCP suggestions，并修复 suggestion 调用里 MCP 分区参数、排名来源与后缀文案等问题。

\textbf{功能说明}：
\begin{itemize}
\item 修复 Markdown 报告渲染检测（\texttt{fix markdown report rendering detection}）
\item 为报告渲染增加 Markdown 标题与表格样式（\texttt{add markdown heading and table styles for report rendering}）
\item 改进报告指标相关展示/错觉问题（\texttt{improve report indicator illusions}）
\item 约束结构化报告指标来源（\texttt{constrain structured report indicator sources}）
\item 用 MCP 建议丰富车辆报告（\texttt{enrich vehicle reports with mcp suggestions}）
\item 修复车辆 suggestion 的 MCP partition 参数（\texttt{fix vehicle suggestion mcp partition arg}）
\item 修复车辆排名来源与 suggestion 后缀（\texttt{fix vehicle rank source and suggestion suffix}）
\item 移除固定车辆 suggestion 后缀（\texttt{remove fixed vehicle suggestion suffix}）
\end{itemize}

\section{可观测性与工程卫生：解析失败内容透出、鉴权 session 兼容、文档与忽略规则}

\textbf{背景}：结构化报告 JSON 解析失败时，若前端只看到空白，不利于排障与用户反馈；auth session schema 变更需兼容旧会话；仓库侧补充忽略规则与文档整理。

\textbf{方案说明}：暴露结构化报告解析失败时的原始/可读内容；修复 auth session schema 兼容性；忽略 Python 缓存与临时文件；清理冗余手册类文档。

\textbf{功能说明}：
\begin{itemize}
\item 前端在结构化报告解析失败时增加失败内容透出，便于排查问题（\texttt{expose structured report parse failure content}）
\item 修复 auth session schema 兼容性（\texttt{fix auth session schema compatibility}）
\item 忽略 Python cache 与临时文件（\texttt{ignore python cache}、\texttt{ignore temporary files}）
\item 文档杂项更新（\texttt{docs: chores and update}）
\item 移除冗余手工文档（\texttt{rm redundant manual docs}）
\end{itemize}

\section{OCR 审核与部署：修正小的架构问题后完成合并与上线准备}

\textbf{背景}：RAG 主链路接入本周没有继续展开，重点是审核并合并 Chessman 完成的 OCR 开发工作。在审核过程中发现了一个小的架构问题，做了修正，再完成后续合并与部署。

\textbf{方案说明}：本周没有继续扩展 RAG 主链路能力，而是先处理 OCR 这条已经在推进的工作线。做法上先完成代码审核，修正其中一个小的架构问题，再把这部分工作合并进主线并完成部署，确保 OCR 能力先稳定落地。

\textbf{功能说明}：
\begin{itemize}
\item 完成了 OCR 相关工作的代码审核，并在合并前修正了一个小的架构问题。
\item 完成了同学 OCR 工作的合并，减少后续主链路接入前的分支分散状态。
\item 完成了 OCR 相关部署，下周可以推进 RAG 知识库增删改查和调用测试。
\end{itemize}

\section{技术债务与限制}

\begin{itemize}
\item \textbf{继承上周未完成：opening/closing 体验仍依赖阶段 prompt}：虽然多步开场/退场能力已接入主链路，但尚无独立的产品化配置层，什么时候该开场、什么时候该退场，仍主要靠任务分类和 prompt 约束。
\item \textbf{继承上周未完成：message sources 下游消费仍未完全统一}：来源信息在 worker、报告和咨询链路里已更一致，但前端展示、埋点和调试面板是否都能稳定消费，仍需继续验证。
\item \textbf{继承上周未完成：\texttt{runtime.ts} 主体拆分已基本完成，但剩余组合逻辑仍需继续减少}：本周已继续拆出 HTTP、Query/OpenAI、Auth、KB、Rule、ToolProvider 与 prompt 模块，当前剩余约一千行组合逻辑，后续仍需继续梳理 import 图、启动路径与潜在循环依赖。
\item \textbf{继承上周未完成：可靠性脚本与检索检查包仍未进入固定测试节奏}：本周提交层面仍未体现 Playwright/可靠性脚本的例行化执行，后续需要补记录或纳入 CI / 日常回归流程。
\item \textbf{RAG 主链路接入}本周没有继续推进；和检索相关的实际工作主要是 OCR 方向的代码审核、合并与部署，主链路接入与效果验证需要明确移到下周继续开展。
\item \textbf{报告 MCP suggestions 与指标来源}已做多轮修正，仍依赖 MCP 返回稳定性与分区语义一致，线上需持续观察误推荐与指标误读。
\item \textbf{结构化解析失败}虽已透出内容，前端展示方式与用户提示文案是否足够友好，仍可继续打磨。
\item 性能/承载评估尚未开展，需要延后完成。
\end{itemize}

\section{下周计划}

\begin{enumerate}
\item 在本周完成 OCR 审核、合并与部署的基础上，启动 RAG 主链路接入与效果验证，并与 \texttt{retrieval-latest-check}、测试服说明等材料对应起来。
\item \textbf{继续统一各报告管线表现}：在驾驶员/单位/线路与车辆报告之间对齐默认行为、错误提示与卡片渲染边界。
\item \textbf{巩固 Router 与专家分流}：在多轮、lookup 失败回路由、续写场景下抽样真实会话，减少误分流。
\item \textbf{补齐或固化端到端测试}：opening/closing、message sources、结构化报告与 Markdown 混排的关键路径。
\item \textbf{继续减少 Runtime 剩余组合逻辑}：在本周拆分基础上，梳理 import 图与启动路径，避免「拆文件但主文件复杂度没有明显下降」。
\item \textbf{驾驶员专家与车辆专家并行演进}：统一工具描述、scoped calls 与负向提示，减少抢答与越界。
\item \textbf{文档与样例}：保持 MCP usable samples、demo 与线上字段对齐，减少联调伪问题。
\item \textbf{性能/承载评估}：尚未开展，需要延后完成。
\end{enumerate}

\end{document}
